\documentclass[10pt,a4paper]{amsart}
\usepackage[margin=19mm]{geometry}
\usepackage{amsmath,amssymb,mathtools}
\usepackage[scr=rsfs,cal=euler]{mathalfa}
\usepackage{xfrac}
\usepackage[unicode,hidelinks,bookmarks=false]{hyperref}
\newcommand{\ie}{i.e.\ }
\newcommand{\cf}{cf.\ }
\newcommand{\resp}{respectively\ }
\newcommand{\Subsection}{\subsection}
\providecommand{\nobreakdash}{\nobreak\hbox{-}}
\setlength{\emergencystretch}{2em}
\multlinegap0pt
\usepackage{bm}
\usepackage{bbm}
\usepackage{dsfont}
\usepackage{xspace}
\usepackage{cancel}
\usepackage{mathtools}
\usepackage{amsthm}
\makeatletter
\@ifundefined{thm}{\newtheorem{thm}{Thm}}{}
\@ifundefined{lem}{\newtheorem{lem}[thm]{Lemma}}{}
\makeatother
\DeclareMathOperator{\Ext}{Ext}
\DeclareMathOperator{\Hom}{Hom}
\DeclareMathOperator{\modu}{mod}
\DeclareMathOperator{\rad}{rad}
\newcommand{\defl}{\twoheadrightarrow}
\newcommand{\infl}{\rightarrowtail}
\newcommand{\mcE}{\mathcal{E}}
\title[Classification of exact structures using the Ziegler spectru]{Classification of exact structures using the Ziegler spectrum}
\author{Julia Sauter}
\date{Source: arXiv:2506.02304v2 (CC BY 4.0)}
\begin{document}
\maketitle

\noindent\emph{Source and licence.} Classification of exact structures using the Ziegler spectrum. Version arXiv:2506.02304v2 (CC BY 4.0), distributed under the Creative Commons Attribution 4.0 International licence, as recorded in the arXiv licence metadata for this exact version and shown on the abstract page https://arxiv.org/abs/2506.02304v2. Full rights notice: licenses/arxiv-2506.02304-CC-BY-4.0.txt.\par\smallskip

\noindent\textit{Editorial scope.} Counted unit: the Lem whose source label is \texttt{radExt} in the article (source0.tex lines 1237--1246, source label \texttt{radExt}), delivered together with the article's own complete original proof of it (source0.tex lines 1248--1285, one \texttt{proof} environment of 38 lines). No mathematics is written, repaired or re-derived: statement and proof are the article's own text line for line. Required context retained uncounted: none. Every definition, display and label of those lines is retained unchanged. Omitted from this package and not redistributed: 1--1236, 1247--1247, 1286--1853. Nothing inside a retained excerpt is omitted or altered: every retained body is delivered exactly as the article's lines are written, so no omission receipt of any kind accompanies this package. Citations: the retained excerpts carry no citation command at all, so no bibliography entry of the article is transcribed and no reference is redistributed. Reference adaptation: none was needed and none was made; every reference inside the delivered unit resolves inside it, so no unresolved reference or citation is produced. Printed slips kept literal: no printed word, symbol, hypothesis or number of the counted statement or of its proof was altered. Layout and class: the article's own class is \texttt{amsart} with packages including amssymb, amscd, fancyhdr, amsmath, amsfonts, stmaryrd, xcolor, hyperref, amsrefs, inputenc, tikz, tikz-cd, todonotes; the wrapper is the lane's pinned \texttt{amsart} editorial wrapper at 10pt on A4, which restates the theorem-like environments the retained text needs and adds the article's own macro definitions that the retained text needs (\texttt{\textbackslash Ext}, \texttt{\textbackslash Hom}, \texttt{\textbackslash defl}, \texttt{\textbackslash infl}, \texttt{\textbackslash mcE}, \texttt{\textbackslash modu}, \texttt{\textbackslash rad}), with extra wrapper packages (bm, bbm, dsfont, xspace, cancel, mathtools). The delivered numbering is the unit's own. Assets: no retained span contains a figure environment, a graphics-inclusion command, a PDF-page inclusion, a table or a TikZ picture; no image file is copied into this package, the package declares no asset dependency and no image file is redistributed with it. Licence: version arXiv:2506.02304v2 of this article is distributed under the Creative Commons Attribution 4.0 International licence, as recorded in the arXiv licence metadata for this exact version and shown on the abstract page https://arxiv.org/abs/2506.02304v2. A full rights notice is delivered at licenses/arxiv-2506.02304-CC-BY-4.0.txt. Characters: every retained excerpt is pure ASCII, so the delivered engine, XeLaTeX with its default Latin Modern text font, reports no missing character.
\begin{lem} \label{radExt} Let $R$ denote a commutative discrete valuation ring   with maximal ideal $P$ and let $n \in\mathbb{N}$.  
We have a functor $\rad^n \colon R\modu \to R\modu$ defined by $\rad^n M = P^n M$. Let $L=[1,n]\subseteq \mathbb{N}$. 
\begin{itemize}
\item[(1)] Then we have $\Ext^1_{\mcE_L} (-,-)=\rad^n\Ext^1_R(-,-)$ and $\Ext^1_{\mcE_L'} (-,-)=\rad^n_R\Ext^1_{\mcE'}(-,-)$
\item[(2)] For every $\mcE_L$-exact sequence $\sigma$ and every object $X$, the sequences $\Ext^1_{\mcE_L} (X,\sigma)$ is right exact.\\
\noindent
For every $\mcE'_L$-exact sequence $\sigma$ and every object $X$, the sequences $\Ext^1_{\mcE'_L} (X,\sigma)$ is right exact.
\end{itemize}    
In particular, $\mcE_L$ and $\mcE'_L$ are hereditary exact. 
\end{lem}

\begin{proof}
\begin{itemize}
\item[(1)] We first describe the $R$-module structure on 
\begin{itemize}
    \item[(a)] $\Ext^1_R (R/P^m, R/P^{\ell})\cong R/P^s$ where $s=\min (m, \ell)$. For $a,b\in \mathbb{N}$ we write $\sigma_{a,b}\in \Ext^1_R(R/P^m, R/P^{\ell})$ for an exact sequence with middle term $R/P^a\oplus R/P^b$ (whenever this exist). In $R/P^s$ we pick $P=(p)$ and we have the following mult. by $p$ \\
    $1 \mapsto p \mapsto p^2 \mapsto \cdots \mapsto p^{s-1} \mapsto 0$
    this corresponds to the following on the Ext-group
    \begin{itemize}
     \item[(a1)] If $s=\ell \leq m$ we have \\
      $\sigma_{0, m+\ell}\mapsto \sigma_{1, m+\ell-1} \mapsto\cdots \mapsto  \sigma_{\ell-1, m+1}  \mapsto 0$\\
     Then we have ${\rm rad}^n\Ext^1_R(R/P^m, R/P^{\ell})\cong R/P^{s-n}$ whenever $n<s$ and zero otherwise. For $n<s=\ell$ it is the image of $p^n$, i.e.   \\
      $\sigma_{n, m+\ell-n}\mapsto \sigma_{n+1, m+\ell-n-1} \mapsto\cdots \mapsto  \sigma_{\ell-1, m+1} $
     \item[(a2)] If $\ell >m=s$ we have \\
     $\sigma_{\ell +m, 0}\mapsto \sigma_{\ell+m-1, 1} \mapsto\cdots \mapsto  \sigma_{\ell+1,m-1} \mapsto 0$\\
     Then we have ${\rm rad}^n\Ext^1_R(R/P^m, R/P^{\ell})\cong R/P^{s-n}$ whenever $n<s$ and zero otherwise. For $n<s=\ell$ it contains the following elements  \\
      $\sigma_{\ell+m-n, n}\mapsto \sigma_{\ell +m-n-1,n+1} \mapsto\cdots \mapsto  \sigma_{\ell+1, m-1} $
    \end{itemize}
    \item[(b) ] $\Ext^1_R(R/P^m, R)\cong R/P^m$ we write $\sigma_a$ for the extension with $R \oplus R/P^a$ as middle term. Then the multiplication by $p$ is given by \\
    $\sigma_{0} \mapsto \sigma_{1} \mapsto \cdots \mapsto \sigma_{m-1} \mapsto 0$\\
    The submodule ${\rm rad}^n\Ext^1_R(R/P^m, R)$ is of course zero is $n\geq m$ and if $n<m$ is given by the following \\
    $\sigma_{n} \mapsto \sigma_{n+1} \mapsto \cdots \mapsto \sigma_{m-1}$
\end{itemize} 
We claim $\Ext^1_{\mcE_{[0,n]}}=\rad^n\Ext^1_R$. First observe that in $\mcE_L$: $R/P^a$ $1\leq a \leq n$ are injectives and they are also projectives. So for $s= \min (\ell ,m) \leq n$ we have $\Ext^1_{\mcE_{[0,n]}}(R/P^m, R/P^{\ell})=0$ and for $m\leq n $ we have $\Ext^1_{\mcE_{[0,n]}}(R/P^m, R)=0$. \\
So we may always assume wlog that $n< \min (\ell,m)$, then proceed by induction over $n$. For $n=1$, a short exact sequence of in $R\modu$ is in $\mcE_{[0,1]}$ iff the indcomposable summands of number of the indec. summands of the outer terms add up to the indec summands of the middle term. This means in case (a) all exact sequence are in this exact structure except $\sigma_{0, m+n}$ and $\sigma_{m+n,0}$, in case (b) all except $\sigma_0$. \\
For $n>1$, it is enough to observe that $\Hom (-, R/P^n)$ is 
\begin{itemize}
\item[(ad a1)] exact on $\sigma_{n, m+\ell -n}$ and not exact on  $\sigma_{n-1, m+\ell -n+1}$
\item[(ad a2)] exact on $\sigma_{\ell+m-n, n}$ and not exact on  $\sigma_{\ell +m -n+1, n-1}$
\item[(ad b)] exact on $\sigma_{n}$ and not exact on  $\sigma_{n-1}$
\end{itemize}
Then the rest follows by induction hypothesis. \\
Now, for $\mcE'_L$ one observes that $\Ext^1_{\mcE'_L}(X,Y)=\Ext^1_{\mcE_L}(X,Y)$ for all $X,Y$ torsion, and for $Y$ free it is zero.\\
As $\Ext^1_{\mcE'}(X,Y)=\Ext^1_R(X,Y)$ for all $X,Y$ torsion and for $Y$ free it is zero. Then the claim $\Ext^1_{\mcE'_L}$ follows from the proof for $\Ext^1_{\mcE_L}$. 
\item[(2)] Taking $n$-th radical of an epimorphism in $R \modu$ is again an epimorphism - as the $n$-th radical can be described as the image of multiplication by $p^n$ (making it also a quotient and not ony a submodule). As $\Ext^1_R(X, \sigma)$ (resp. $\Ext^1_R(\sigma, X)$) is right exact for all exact sequences $\sigma$, we conclude that $\rad^n \Ext^1_R(X, f)$ (resp. $\rad^n \Ext^1_R(g, X)$) are epimorphisms for $f$ an epimorphism and $g$ a monomorphism. \\
Of  course, on general exact sequence $\rad^n$ is not a middle-exact functor (but this follows from (1) since for $\sigma \in \mcE_L$ the sequences $\Ext^1_{\mcE_L}(X,\sigma)$ and $\Ext^1_{\mcE_L}(\sigma ,X)$ are middle exact). In particular, the right exactness claim follows. \\
For $\Ext^1_{\mcE'_L}(X,\sigma) $ we can restrict to $X$ indecomposable torsion and as $\sigma$ in $\mcE'$, it fits into an exact sequence of ses $\sigma_{tor}\infl \sigma \defl \sigma_{free}$ with $\sigma_{tor}$ the torsion part and $\sigma_{free}$ the free part, and we conclude that $\Ext^1_{\mcE'_L}(X,\sigma)= \Ext^1_{\mcE'_L} (X, \sigma_{tor}) =\Ext^1_{\mcE_L} (X, \sigma_{tor})$ is right exact.   
\end{itemize}
\end{proof}

\end{document}
